\subsection{应用：拉普拉斯变换}

在本号中，M 表示一个交换幺半群，其运算用加法记号表示，并赋予局部紧空间的拓扑，使得映射 \((m, m') \mapsto m + m'\) 从 \(M \times M\) 到 M 连续。M 的单位元记为 0。称 M 的特征为 M 上的每个有界连续复函数 \(\chi\)，满足以下关系：

\[
\chi (m + m ^ {\prime}) = \chi (m) \cdot \chi (m ^ {\prime}), \quad \chi (0) = 1, \quad | \chi (m) | \leqslant 1\tag{7}
\]

对 \(m, m'\)（属于 M），若 \(\chi\) 和 \(\chi'\) 是特征，则 \(\chi\chi'\) 也是特征。M 的特征全体构成一个幺半群，记为 X；赋予 X 紧收敛拓扑，使从 X × X 到 X 的映射 \((\chi, \chi') \mapsto \chi\chi'\) 连续。X 的单位元是常函数 1。

对 M 上的每个有界复测度 \(\mu\)，称 \(\mu\) 的拉普拉斯变换为定义在 X 上的函数 \(\mathcal{L}\mu\)，由下式给出：

\[
(\mathscr {L} \mu) (\chi) = \int_ {\mathrm{M}} \chi (m) d \mu (m).\tag{8}
\]

由 GT, X, §3, No.~4 的定理 3，M×X 到 C 的映射 \((m,\chi)\mapsto\chi(m)\) 连续且有界。于是第 6 号的命题 13 的推论给出以下结果：

命题 14. --- 对 M 上的每个有界复测度 \(\mu\)，函数 \(\mathcal{L}\mu\) 在 X 上连续且有界。若 \(\mathcal{M}_{+}^{b}(M)\) 赋予窄拓扑、\(\mathcal{C}^{b}(X; \mathbf{C})\) 赋予紧收敛拓扑，则映射 \(\mu \mapsto \mathcal{L}\mu\) 从 \(\mathcal{M}_{+}^{b}(M)\) 到 \(\mathcal{C}^{b}(X; \mathbf{C})\) 连续。

M 上在无穷远处趋于 0 的特征所成的集合记为 \(X_{0}\)；该集合对乘法稳定。称 \(^{(1)}\) S 是 X 的满子幺半群，如果 S 对映射 \(\chi\mapsto\overline{\chi}\) 稳定，\(S\cap X_{0}\) 分离 M 的点（GT, X, §4, No.~1, 定义 1），并且对任意 \(m\in M\)，存在 \(\chi\) 中的元素 \(S\cap X_{0}\)，使得 \(\chi(m)\neq0\)。

再假设 M 是非紧交换群。设 f 是 M 上在无穷远处趋于 0 的函数；则 M 上的函数 \(x \mapsto f(x)f(-x)\) 也在无穷远处趋于 0，而 M 的每个特征 \(\chi\) 都满足 \(\chi(x)\chi(-x) = \chi(0) = 1\)。因此 \(X_{0}\) 为空，且 X 不包含任何满子幺半群。所以，下面的定理 3 不适用于非紧的局部紧群。

定理 3. --- 设 S 是 X 的满子幺半群。

\begin{enumerate}
\def\labelenumi{\alph{enumi})}
\item
  若 \(\mu\) 和 \(\mu'\) 是 \(\mathbf{M}\) 上的两个有界复测度，且 \(\mathcal{L}\mu\) 和 \(\mathcal{L}\mu'\) 在 \(\mathrm{S} \cap \mathrm{X}_0\) 上的限制相同，则 \(\mu = \mu'\)。
\item
  设 \(\mathfrak{F}\) 是 \(\mathcal{M}_{+}^{b}(\mathbf{M})\) 上的滤子，并且 \(\mathcal{L}\lambda(s)\) 有极限 \(\Phi(s) \in \mathbf{C}\)，关于 \(\mathfrak{F}\)，对每个 \(s \in \mathbf{S}\) 均如此。则滤子 \(\mathfrak{F}\) 淡收敛到有界正测度 \(\mu\)，并且 \(\Phi(s) = \mathcal{L}\mu(s)\) 对所有 \(s \in \mathbf{S} \cap \mathbf{X}_0\) 成立。
\item
  在 \(b\) 的假设下，再假设 \(S \cap X_0\) 的闭包包含 1，且函数 \(\Phi\) 在 \(S\) 上于点 1 连续。则 \(\mathfrak{F}\) 窄收敛到 \(\mu\)，并且 \(\Phi(s) = \mathcal{L}\mu(s)\) 对所有 \(s \in S\) 成立。
\end{enumerate}

记 \(S \cap X_0\) 为 M 上在无穷远处趋于 0 的连续复函数代数，记 A 为由 \(\mathfrak{A}\) 生成的 E 的线性子空间；则 \(E\) 是 E 的子代数，并且对映射 \(f \mapsto \overline{f}\) 稳定；由于 \(S\) 是 \(X\) 的满子幺半群，GT, \(X\)，§4, No.~4 的命题 7 的推论 2 蕴含 \(\mathfrak{A}\) 在 \(E\) 中稠密。

证明 \(a\)：根据假设，对每个 \(\mu(f) = \mu'(f)\) 有 \(f \in \mathfrak{A}\)；由于 \(\mu\) 和 \(\mu'\) 是 \(\mathbf{E}\) 上的连续线性形式，这就蕴含对 \(\mu(f) = \mu'(f)\) 有 \(f \in \mathbf{E}\)，特别地，对每个具有紧支集的连续函数 \(f\) 都成立，从而 \(\mu = \mu'\)。

现在置于 \(b\) 的假设下。数 \(\Phi(1) = \lim_{\lambda, \mathfrak{F}} \lambda(1)\) 为实数且为正；取给定的实数 \(a > \Phi(1)\)；由于

由 \(\| \lambda \| = \mathcal{L}\lambda (1)\) 对 \(\lambda \in \mathcal{M}_{+}^{b}(\mathrm{M})\) 成立，关系 \(\lim_{\lambda ,\mathfrak{F}}\mathcal{L}\lambda (1) = \Phi (1)\) 蕴含由满足 \(\lambda \in \mathcal{M}_{+}^{b}(\mathrm{M})\) 且 \(\| \lambda \| \leqslant a\) 的测度所成的集合 H 属于 \(\mathfrak{F}\)。由于 \(\mathcal{M}^b (\mathrm{M};\mathbf{C})\) 可与赋范空间 E 的对偶空间相识别（第 III 章 §1, No.~8 \& §1, No.~2, 命题 3），空间 H 关于拓扑 \(\sigma (\mathcal{M}^b (\mathrm{M};\mathbf{C}),\mathrm{E})\) 是紧的。另一方面（TVS, III, §3, No.~4, 命题 5），该拓扑在 H 上与 E 的任意总子集上的逐点收敛拓扑一致。特别地，由于 \(\mathfrak{A}\) 在 E 中稠密，且具有紧支集的连续函数空间也在 E 中稠密（第 III 章 §1, No.~2, 命题 3），所以 \(\mathrm{S}\cap \mathrm{X}_0\) 上的逐点收敛拓扑在 H 上与淡拓扑一致，H 关于此拓扑是紧的。立即可知，\(\mathfrak{F}\) 淡收敛到某个测度 \(\mu \in \mathrm{H}\)，并且 \(\mathcal{L}\mu (s) = \lim_{\lambda ,\mathfrak{F}}\mathcal{L}\lambda (s)\) 对所有 \(s\in \mathrm{S}\cap \mathrm{X}_0\) 成立。

最后处理 \(c\)。由于函数 \(\Phi\) 和 \(\mathcal{L}\mu\) 在 \(1 \in S\) 处连续，在 \(S \cap X_0\) 上相等，且 \(S \cap X_0\) 的闭包包含 1，有 \(\Phi(1) = \mathcal{L}\mu(1)\)。换言之，\(\lim_{\lambda, \mathfrak{F}} \lambda(1) = \mu(1)\)。第 3 号命题 9 遂说明 \(\mu\) 是滤子 \(\mathfrak{F}\) 的窄极限。\(S\) 的每个元素都是 \(M\) 上的有界连续函数，因此对所有 \(\Phi(s) = \lim_{\lambda, \mathfrak{F}} \lambda(s) = \mu(s) = \mathcal{L}\mu(s)\) 有 \(s \in S\)。

推论。 --- 设 S 是 X 的满子幺半群，且 \(S \cap X_{0}\) 的闭包包含 1。令 L 为 \(\mathcal{C}^{b}(S; \mathbf{C})\) 的子集，由 \(\lambda \in \mathcal{M}_{+}^{b}(M)\) 的测度的拉普拉斯变换在 S 上的限制组成。

\begin{enumerate}
\def\labelenumi{\alph{enumi})}
\item
  集合 L 在赋予逐点收敛拓扑的空间 \(\mathcal{C}^{b}(\mathrm{S};\mathbf{C})\) 中是闭的。
\item
映射 \(\lambda \mapsto (\mathcal{L}\lambda)_\mathrm{S}\) 将 \(\mathcal{M}_+^b (\mathbf{M})\) 同胚映到 \(\mathbf{L}\)；若 \(\mathcal{M}_+^b (\mathbf{M})\) 赋予窄拓扑、\(\mathbf{L}\) 赋予逐点收敛拓扑，则该映射是同胚。
\item
  L 上的逐点收敛拓扑与紧收敛拓扑一致。
\end{enumerate}

断言 \(a)\) 和 \(b)\) 是定理 3 的直接推论；断言 \(c)\) 由 \(b)\) 和命题 14 得出，因为紧收敛拓扑比逐点收敛拓扑细。

必须注意，L 在赋予逐点收敛拓扑的 S 上所有有界复函数所成的集合中并不闭。例如，采用下面例 2 的记号（M = R+，S 视同 R+）。测度 \(\varepsilon_{n}\)（\(n \in N\)）的拉普拉斯变换是 R+ 上的函数 \(t \mapsto e^{-nt}\)；当 n 趋于 \(+\infty\) 时，这些函数逐点收敛到在 t = 0 时等于 1、在 \(t \neq 0\) 时等于 0 的函数，而该函数不属于 L。

例 1)。 --- 取 M 为正整数集合 N，赋予加法运算和离散拓扑。令 D 为 C 的单位圆盘（绝对值 ≤ 1 的复数集合），赋予 C 诱导的拓扑和乘法诱导的运算。对每个 z ∈ D，记 f(z) 为 N 的特征 n ↦ z\^{}n。对 N 的每个特征 χ，记 g(χ) 为复数 χ(1) ∈ D。立即验证 f 和 g 互为 D 与 X 之间的同胚，这使我们从现在起可以视 X 与 D 为同一对象。在无穷远处趋于 0 的特征于是可视同绝对值 \textless{} 1 的复数所成的集合 D₀。最后，R 的区间 {[}0, 1{]} 是 D 的满子幺半群，且 1 属于 {]}0, 1{]} ∩ D₀ = {]}0, 1{[} 的闭包。

每个测度 \(\mu\) 定义在 \(\mathbf{N}\) 上，都能唯一写成 \(\mu = \sum_{n\in \mathbf{N}}u_n\cdot \varepsilon^n\) 的形式，并且 \(\mu\) 有界当且仅当 \(\sum_{n}|u_n| < +\infty\)；此时 \(\mathcal{L}\mu (z) = \sum_{n\in \mathbf{N}}u_nz^n\)，对 \(z\in \mathbf{D}\) 成立。此函数在 \(\mathbf{D}\) 上连续；通常称其为可求和序列 \((u_n)_{n\in \mathbf{N}}\) 的生成函数。用这种语言表述定理 3，并考虑第 3 号的命题 9，得到以下结果：

命题 15. --- 设 A 是赋予滤子 F 的集合。对每个 \(\alpha\in A\)，设 \((u_{\alpha,n})_{n\in\mathbf{N}}\) 是正数的可求和序列，令 \(\Phi_{\alpha}\) 为定义在 R 的区间 {[}0,1{]} 上的函数，由 \(\Phi_{\alpha}(x)=\sum_{n\in\mathbf{N}}u_{\alpha,n}x^{n}\) 给出。要使存在正数的可求和序列 \((u_{n})_{n\in\mathbf{N}}\)，使得

\[
\lim _ {\alpha , \mathfrak {F}} u _ {\alpha , n} = u _ {n} \text {   for   all   } n, \quad \lim _ {\alpha , \mathfrak {F}} \sum_ {n \in \mathbf {N}} u _ {\alpha , n} = \sum_ {n \in \mathbf {N}} u _ {n},
\]

必要且充分的条件是：\(\Phi_{\alpha}\) 在 \(]0,1]\) 上关于 \(\mathfrak{F}\) 逐点收敛到在点 1 连续的函数 \(\Phi\)。此时 \(\Phi(x) = \sum_{n \in \mathbf{N}} u_n x^n\)，对所有 \(x \in ]0,1]\) 成立。

取 M 为幺半群 \(N^{n}\)，其中 n 表示 \textgreater1 的整数；于是可将空间 X 视同 \(D^{n}\)，并可选取 \([0,1]^{n}\) 作为满子幺半群。本情形下定理 3 的转述留给读者。

例 2)。 --- 取 M 为集合 \(R_{+}\)，赋予加法运算和通常拓扑。令 P 为实部为正的复数 z 所成的集合，赋予 C 诱导的拓扑和 C 中加法诱导的运算。对每个 \(p \in P\)，记 \(f(p)\) 为 \(x \mapsto e^{-px}\) 的特征 \(R_{+}\)；容易验证，f 是 P 的拓扑幺半群结构到 X 的拓扑幺半群结构的同构；我们通过 f 将 X 视同 P。显然，\(R_{+}\) 是 P 的满子幺半群，且定理 3 给出以下结果：

命题 16. --- 设 A 是赋予滤子 \(\mathfrak{F}\) 的集合。对每个 \(\alpha \in \mathbf{A}\)，设 \(\mu_{\alpha}\) 是 \(\mathbf{R}_{+}\) 上的有界正测度，并令 \(\Phi_{\alpha}\) 为定义在 \(\mathbf{R}_{+}\) 上的函数，由 \(\Phi_{\alpha}(p) = \int_{0}^{+\infty} e^{-px} d\mu_{\alpha}(x)\) 给出。要使映射 \(\alpha \mapsto \mu_{\alpha}\) 关于 \(\mathfrak{F}\) 窄收敛到有界正测度 \(\mu\)，必要且充分的条件是：\(\Phi_{\alpha}\) 在 \(\mathbf{R}_{+}\) 上关于 \(\mathfrak{F}\) 逐点收敛到在点 0 连续的函数 \(\Phi\)。此时 \(\Phi(p) = \int_{0}^{+\infty} e^{-px} d\mu(x)\)，对所有 \(p \in \mathbf{R}_{+}\) 成立。

可加幺半群 \(\mathbf{R}_{+}^{n}\) 也有类似结果（n 是 \textgreater1 的整数）；本情形下定理 3 的转述留给读者。

